\documentclass[12pt]{article}
\usepackage[utf8]{inputenc}
\usepackage[T1]{fontenc}
\usepackage{hyperref}
\title{Tzeltal Language Resources for K-12 Curriculum Development}
\author{Compiled by Claude Sonnet 5 (claude-sonnet-5) for Andr\'as Kornai}
\date{}
\begin{document}
\maketitle
\section{Introduction}
Tzeltal (ISO 639-3 tzh) is a Mayan language of Chiapas, Mexico, with about 450,000 speakers. Search found mostly older materials and library catalogue records. It found no current, comprehensive course.
\section{Primary Educational Resources}
\begin{itemize}
\item A Tzeltal reader from the ``Bats'i K'op'' series, first cycle, appears in a library catalogue. \url{https://opac.biblio.iteso.mx/vufind/Record/000499546/Channels}
\item A 1974 grammar of Spanish with explanations in Tzeltal (Ocosingo variety), held by El Colegio de M\'exico. \url{https://repositorio.colmex.mx/concern/books/7d278v896?locale=en}
\item A 1954 Tzeltal primer (Lacand\'on variety), 42 pages. The catalogue note says its language is no longer current. \url{https://repositorio.colmex.mx/concern/books/c534fq90b?locale=en}
\end{itemize}
\section{Linguistic and Reference Materials}
\begin{itemize}
\item Diccionario de la lengua tzeltal (1888), catalogued by the BVFE. \url{https://www.bvfe.es/en/directorio-bibliografico-diccionarios-vocabularios-glosarios-tratados-y-obras-lexicografia/17540-diccionario-de-la-lengua-tzeltal/imprimir.html}
\item Tzeltal dictionary PDF (second edition) hosted by Arqueolog\'ia Mexicana. I could not confirm its contents. \url{https://arqueologiamexicana.mx/sites/default/files/banco_imagenes/tzh_diccionario_ed2.pdf}
\item Another El Colegio de M\'exico holding (Oxchuc phrasebook or related). \url{https://repositorio.colmex.mx/concern/books/c534fq89k?locale=en}
\item Another El Colegio de M\'exico holding. \url{https://repositorio.colmex.mx/concern/books/kd17cw054?locale=en}
\item Lexicity Mayan page, which lists a Tzeltal grammar and bibliography. \url{https://oderalon.net/wiki/index.php/Lexicity/Mayan}
\end{itemize}
\section{Religious and Community Resources}
None found in this search.
\section{Digital and Audio Resources}
A Smithsonian repository item turned up in the results. I could not confirm what it contains. \url{https://repository.si.edu/items/318b5c11-94f8-46c6-a6ae-767cf580d5a5/full}
\section{Curriculum Development Context}
Existing materials are mostly dated and dialect-specific (Bachaj\'on, Ocosingo, Oxchuc, Lacand\'on). A curriculum would need new, current materials. INALI and SEP bilingual-education materials are worth checking, but I did not find them in these results.
\end{document}